% !TeX root = proyecto.tex
% RESUMEN / ABSTRACT

\clearpage
\thispagestyle{empty}

\begin{center}
    {\Large\bfseries\color{titleColor} Resumen}
\end{center}

\vspace{0.8cm}

\noindent
El presente Trabajo Terminal (no.\ 2026-A135) constituye la segunda fase de un proyecto de
investigación aplicada orientado a mitigar la invisibilidad estadística de las Niñas, Niños y
Adolescentes (NNA) en situación de orfandad como consecuencia directa de un feminicidio en México.
Ante la ausencia de un registro nacional unificado y la inviabilidad operativa del monitoreo
manual del ecosistema periodístico digital, se diseñó y desarrolló un \textbf{sistema inteligente}
que automatiza el ciclo completo de descubrimiento informativo: recolección, deduplicación,
clasificación semántica e investigación profunda de casos.

\vspace{0.5cm}

\noindent
El módulo de recolección implementa una arquitectura híbrida \textbf{RSS + StealthSession} que
monitorea 45 fuentes periodísticas nacionales con cobertura complementaria mediante emulación de
\textit{fingerprint} TLS, eludiendo los \textit{Web Application Firewalls} que bloqueaban el
acceso directo. La clasificación semántica se basa en \textbf{BETO}
(\textit{bert-base-spanish-wwm-cased}), un modelo Transformer bidireccional de 110 millones de
parámetros preentrenado en español, integrado en una \textbf{arquitectura de clasificación
escalonada} de cuatro capas que combina reglas simbólicas de dominio con la inferencia neuronal
para reducir el solapamiento semántico propio del corpus de noticias sobre violencia de género.

\vspace{0.5cm}

\noindent
La evaluación del clasificador se realizó sobre un conjunto de prueba etiquetado de forma
independiente por dos anotadores, cuyo acuerdo inter-anotador fue cuantificado mediante el
coeficiente $\kappa$ de Cohen. Los resultados muestran una reducción significativa en la tasa de
falsos positivos respecto al sistema léxico del TT1, ampliando la cobertura geográfica del
monitoreo y reduciendo la latencia en la detección de nuevos casos. La información clasificada
se persiste en una base de datos \textbf{PostgreSQL} con búsqueda de texto completo en español
y es accesible mediante una interfaz web que incorpora el módulo \texttt{DeepInvestigator}, el
cual genera fichas de caso estructuradas con datos accionables para organizaciones civiles e
instituciones gubernamentales, contribuyendo a la materialización del principio del
\textit{interés superior de la niñez} consagrado en la Ley General de los Derechos de Niñas,
Niños y Adolescentes (LGDNNA).

\vspace{1.2cm}

\noindent
{\bfseries\color{sectionColor} Palabras clave:}

\vspace{0.4cm}

\noindent
Niñas, Niños y Adolescentes (NNA);
feminicidio;
orfandad;
procesamiento de lenguaje natural (PLN);
clasificación escalonada neuro-simbólica;
BETO;
web scraping;
Kappa de Cohen;
microservicios;
derechos de la infancia;
México.

\clearpage
